\noindent\hspace*{\parindent}I begin with the (fictitious) assumption that a society has been formed by the coming together of a very large number of married couples of the most varied ages. Each year a number of marriages will be dissolved by the death of one or the other spouse: let this loss be replaced by a definite number of new married couples joining each year, enough for the size of the society as a whole to remain unchanged. Since the society is assumed to be very large, a proportionately very small part of the widows arising through the husband's death in the course of a year will themselves die during that same year, thus modifying somewhat the number of widows actually present at the year's end. By the end of the second year approximately as many new widows will have been added, but some of the widows remaining from the first year will already have died, so that at the end of the second year the number of widows will not be quite twice as great as at the end of the first. At the end of the third year approximately as many new widows have again been added; but the number present at the beginning of the third year will have suffered a loss almost twice as great as the balance of the second year had shown. One sees that in this way the number of widows does indeed grow continually, but ever more slowly, until it has finally become so great that the annual loss through the widows' deaths balances the increase through the deaths of husbands in the society: the steady state will then set in, the number of widows having reached its maximum and remaining unchanged from then on.

If to the above assumptions we add that both the number of married couples
\SourceLeaf{146}
newly joining each year and the proportion in which the various ages are mixed in such a group remain the same from year to year, and likewise that deaths keep exact pace with the mortality tables, then that steady state will deserve its name in the strictest sense; the number of married couples in the society will remain entirely unchanged, as will the number of widows once it has reached its highest value. It is further clear that, if one imagines a second similar society in which the new entrants are mixed in exactly the same proportion as in the first, but which comprises twice as many married couples, the highest number of widows will also be twice as great; or, to put it generally, the ratio of existing marriages to the number of widows after the steady state has set in will depend not on the size of the association (provided only that it is regarded as remaining constant), but solely on the proportion in which the various ages of the new entrants are mixed; and, if the latter proportion were completely known, the former could be determined a priori by calculation. Moreover, the steady state will set in, if not earlier, at any rate when the original participants have all died, which, if the most extreme cases are to be considered, might possibly be delayed until 80 years after the society first came together. One may, however, assume that after only 45—50 years the actual state will have come very close to the steady state\footnote{\textsc{Kritter}, in his opinion, is already content with 40 years.}. To make the ideas more definite, I shall give certain numbers by way of example; since they are intended only to illustrate the otherwise abstract exposition, however, they make no claim to perfectly exact appropriateness. Let the society consist of 2600 married couples, to which 130 new ones are added annually, while on average just as many leave, namely 70 through the husband's death and 60 through the wife's. Of the 70 widows arising during a year, one dies in that same year, so that at the close of the first year there are 69 widows. Of these, a further 2 have died by the end of the second year, but 69 new widows have again been added, so that 136 are present in all. Of these, 4 die during the third year; and since the newly added widows again number 69, the total at the close of the third year is 201. Progressing in this way ever more slowly, the number of widows may be, after 10 years, 600; after 20 years, 1000; after 30 years, 1200; after 40 years, 1280; after 45 years, 1294; and after 50 years, 1298; a few more are added later, producing the actual maximum of 1300. Here, then, we should have the ratio of existing marriages to the highest number of widows as two to one.

In a real society, where the required conditions hold not with their strict exactness but only on average, things naturally cannot proceed as regularly as in the fictitious one. In the former, not exactly 130 new married couples will join annually, but somewhat more in one year and somewhat fewer in another. Likewise, the various ages of the new entrants will be mixed somewhat differently in one year than in another, or than was assumed in the ideal society; for example, the average age of the men, or that of the women, or the average difference between the two will sometimes be somewhat greater, sometimes somewhat smaller. Finally, deaths in a given number of persons will not follow the mortality tables \textit{exactly}; in one year these will indicate somewhat too few, in another somewhat too many. The result of all this will be that in the real society the number of widows will at one time be somewhat greater than in the ideal one, at another somewhat smaller. A true steady state \textit{in the strictest sense} will never set in in the former, but rather a fluctuation about a mean
\SourceLeaf{147}
state. Thus, at the time when the number of widows in the ideal society would have reached the maximum of 1300 entirely or almost entirely, in the real one there will be fluctuations up and down above and below this number. There may be 1350 at one time, perhaps even 1360, and at another 1250, perhaps only 1240. But it is quite impossible here to set sharp limits and to assert definitely of any number, whichever one may choose, that it could still be reached but the next higher one could no longer be reached. The calculus of probabilities, in its present state, teaches how to reduce such fluctuations to a definite measure, derived, however, not from the extreme but from consideration of all intermediate levels and their relative probability. This is of course by no means the place to pursue this theory further; moreover, its applicability to concrete cases depends on either having mathematically precise knowledge of the determining elements (as is usually the case with games of chance, for example), or having a sufficiently rich body of experience available. Neither holds for subjects such as widows' funds, so that nothing can be done a priori towards such a measurement. But the theory does teach this: to infer from the fluctuations which have occurred in one case those which must with equal justification be expected in another case qualitatively the same but quantitatively different. In their absolute size the fluctuations will be greater in a large society than in a small one; in their relative size the reverse will hold. Thus, if, for example, a society of the above kind had at one time counted up to 60 widows more, at another up to 60 fewer, than the average of 1300, but in such a way that the number 1360, or one close to it, had occurred with some persistence, and likewise the number 1240 or one near it, one would be entitled to infer that in another society, which under otherwise similar circumstances counts on average only 26 married couples, the normal number of widows could just as readily appear increased or diminished by 6. Or, put somewhat differently: just as the number of widows in the large society can fluctuate up and down between 1360 and 1240, it can do so in the small one between 7 and 19; the \textit{absolute} fluctuation, if the size of the large society is 100 times greater than the size of the small one, is ten times as great in the large as in the small; with the \textit{relative} fluctuation, which in the large society (according to the above numbers) would amount to \(4\tfrac8{13}\) percent and in the small to \(46\tfrac2{13}\) percent, the relation is exactly reversed.

So much for the matter itself. As for the name, for the numbers in the examples 1300 and 13 I have just used the designation ‘normal number of widows’ instead of the lengthy circumlocution ‘average number of widows after the steady state has set in’, which will readily be admitted to be not inappropriate; but to my knowledge it is not \textit{usual}. Rather, it is quite common to denote that concept simply by \textit{highest number of widows}, entirely disregarding the irregular fluctuations and considering only the steady state in general as opposed to the preceding time. For anyone with some knowledge of the subject a misunderstanding will not readily be possible here; in any case, this much is certain, and immediately recognisable from the opinion if one reads it with only some attention: \textsc{Kritter} understood the highest number of widows in \textit{this} sense, and only in this sense. P., on the other hand, had something quite different in mind, namely the number beyond which the number of widows should never be able to go with the highest probability. Confusing the two numbers is much the same as confusing the average price of a commodity with its highest price. The objection made above on p. [140] under 3 has thereby been sufficiently proved. But a question about the highest number of widows in P.'s sense of the word is one to which no definite answer can be given at all.
